\section*{Chapitre \ref{ch:rc:short-rate-models} --- Modèles de taux court}

\begin{solution}{exo:rc:short-rate-models:1}
$\sigma B(\tau)/\tau$~: à dix ans $80\times(1-e^{-0.3})/0.3 = 69.1$ points de base~; à
trente ans $80\times(1-e^{-0.9})/0.9 = 52.7$.
\end{solution}

\begin{solution}{exo:rc:short-rate-models:2}
Ses prix d'obligations dépendent de trois constantes ($\kappa$, $\bar x$, $\sigma$) et du taux court
d'aujourd'hui, une famille à trois paramètres de courbes qui ne peut pas épouser une courbe quelconque.
Hull--White fait dépendre du temps le niveau de la dérive, $\vartheta(t)$, choisi de sorte
que chaque facteur d'actualisation soit reproduit à l'identique~; sous forme gaussienne, il vaut $r_t =
f(0,t)+x_t$.
\end{solution}

\begin{solution}{exo:rc:short-rate-models:3}
$\bar x-\sigma^2/(2\kappa^2) = 3~\%-0.0001/0.02 = 2.50~\%$.
\end{solution}

\begin{solution}{exo:rc:short-rate-models:4}
$j_{\max} = \lceil0.184/(0.03/12)\rceil = 74$~: 149 nœuds. Avec des pas
hebdomadaires $\lceil0.184/(0.03/52)\rceil = 319$~: 639 nœuds. Un $\kappa$ petit produit des arbres larges~; la
largeur n'est atteinte qu'après $j_{\max}$ pas, donc les arbres courts restent plus étroits.
\end{solution}

\begin{solution}{exo:rc:short-rate-models:5}
Il faut que chaque prix d'obligation à l'échéance soit une fonction monotone d'une seule
variable d'état, de sorte qu'une valeur critique $x^\star$ découpe la région d'exercice. Avec
deux facteurs, la frontière d'exercice est une courbe dans le plan $(x,y)$, et la
région d'exercice de l'obligation à coupon ne se décompose pas en régions
de zéro-coupon à frontière commune.
\end{solution}

\begin{solution}{exo:rc:short-rate-models:6}
Plus pentue~: un retour à la moyenne plus fort amortit davantage la volatilité des taux longs, si bien
que, pour reproduire les échéances ultérieures à partir de swaps plus courts (dont la volatilité
que le marché maintient élevée), $\sigma(t)$ doit croître plus vite~; avec un $\kappa$ plus grand, l'ajustement peut même exiger
de fortes volatilités tardives. $\kappa$ et $\sigma(t)$ se compensent~: la structure par terme de la
volatilité est partagée entre les deux.
\end{solution}

\begin{solution}{exo:rc:short-rate-models:7}
Erreurs par rapport à 0.026018~: $+1.0\times10^{-4}$ (mensuel), $+6.1\times10^{-5}$
(quinzenal), $+2.6\times10^{-5}$ (hebdomadaire), $-0.2\times10^{-5}$ (demi-hebdomadaire). Diviser le
pas par deux réduit à peu près l'erreur de moitié jusqu'à ce qu'elle change de signe~: ordre un, comme
attendu pour un arbre dont le payoff présente une cassure au strike.
\end{solution}

\begin{solution}{exo:rc:short-rate-models:8}
Recaler les swaptions co-terminales fixe la volatilité de chaque taux de swap~; une \omterm{def:rc:convexity-adjustments-and-constant-maturity-products:spread}{option sur spread CMS}
a aussi besoin de la corrélation entre le taux à dix ans et le taux à deux ans,
qui vaut un dans tout modèle à un facteur~: la volatilité du spread s'effondre et
l'option est fortement sous-évaluée. Utiliser un modèle à deux facteurs, ou une copule des
smiles des deux taux, calibrée sur une corrélation.
\end{solution}

\begin{solution}{pb:rc:short-rate-models:1}
\textbf{1.} Les co-terminales jusqu'à dix ans~: ce sont les swaptions européennes dans lesquelles la
swaption bermudéenne peut être exercée, et sa valeur se situe entre la plus grande d'entre elles et un prix
qui dépend de leur dynamique conjointe.
\textbf{2.} 60.8 à 91.0 points de base, en hausse avec l'échéance.
\textbf{3.} Que la volatilité des taux attendue plus tard (pour des swaps restants plus courts) est
plus élevée que ne le donnerait une réversion à la moyenne avec un $\sigma$ constant.
\textbf{4.} Le modèle le recale par construction grâce à $\vartheta(t)$.
\textbf{5.} Il est mal identifié à partir d'une seule colonne de co-terminales~; il se compense
avec $\sigma(t)$. Les desks le fixent par convention ou à partir du lien entre
les volatilités des co-terminales et des autres swaptions, et le gardent stable.
\textbf{6.} 86.0 points de base.
\textbf{7.} 76.0 à 76.7 points de base.
\textbf{8.} $+15.3$ points de base à un an, $-14.2$ à neuf ans.
\textbf{9.} Les exercices précoces sont surévalués (les échéances courtes du modèle sont trop
volatiles), les tardifs sous-évalués.
\textbf{10.} Avec $\kappa=3~\%$ et $\sigma$ constant, les volatilités des co-terminales sont quasiment plates selon
l'échéance, tandis que celle du marché augmente de 30 points de base.
\textbf{11.} 68.7 points de base la première année, 116.3 au maximum (années sept
à huit).
\textbf{12.} La dernière co-terminale (neuf ans dans un an) a une volatilité de 91.0, à peine
supérieure à celle de huit ans~; pour la recaler, les niveaux antérieurs étant
maintenus, le dernier intervalle en demande moins.
\textbf{13.} Exactement les co-terminales~; les autres swaptions seulement dans la mesure où
$\kappa$ et le $\sigma(t)$ ajusté les reproduisent par hasard.
\textbf{14.} Qu'il reste au dernier niveau, 104.2 points de base.
\textbf{15.} Des paramètres bruyants et des sauts de P\&L d'un jour à l'autre~; des ratios de
couverture qui varient avec la calibration plutôt qu'avec le marché. Lisser ou régulariser, et
surveiller les paramètres.
\textbf{16.} Hull--White avec $\kappa$ fixé à 3~\% et un $\sigma(t)$ constant par morceaux sur les
échéances des co-terminales, calibré chaque jour sur les swaptions co-terminales à la monnaie~;
ses limites~: un seul facteur, smile plat, pas d'ajustement aux autres swaptions.
\textbf{17.} Montrer les erreurs d'ajustement des co-terminales avec une volatilité et avec neuf volatilités, la
stabilité de $\sigma(t)$ sur les mois passés, et le prix de la swaption bermudéenne dans les deux cas,
face à un modèle de référence (benchmark), soit un modèle de marché, chapitre~8.
\textbf{18.} Une corrélation imparfaite entre les taux exercés, ce qui abaisse la
valeur de switch entre les dates d'exercice~; le coût est un pricing plus lent, davantage de
paramètres et une corrélation à marquer.
\textbf{19.} Résultat à retenir~: \emph{un seul sigma ne suffit pas}~: la meilleure volatilité constante
rate les co-terminales de $+15.3$ points de base à un an et de $-14.2$ à neuf ans~;
le $\sigma(t)$ ajusté va de 68.7 à 116.3 points de base.
\textbf{20.} Recaler exactement la courbe et les volatilités des options que le produit
peut devenir.
\end{solution}

\begin{solution}{iq:rc:short-rate-models:1}
Il ramène le taux court vers sa trajectoire moyenne, si bien que les chocs s'éteignent et que
les taux longs bougent moins que les taux courts~: il fixe la façon dont la volatilité décroît avec
la maturité ($\sigma B(\tau)/\tau$), et le degré de corrélation des exercices d'une swaption bermudéenne.
Un trader le lit comme la pente de la volatilité du modèle selon les maturités.
\iqlookfor{retour à la moyenne comme paramètre de forme de la volatilité.}
\end{solution}

\begin{solution}{iq:rc:short-rate-models:2}
$r$ est un processus d'Ornstein--Uhlenbeck~; $\int_t^Tr_s\,ds$ sachant $r_t$ suit une loi normale de moyenne $\bar
x\tau+(r_t-\bar x)B(\tau)$ et de variance $\frac{\sigma^2}{\kappa^2}(\tau-B-\frac\kappa2B^2)$~; donc
$P = \exp(-\text{moyenne}+\text{variance}/2)$, une fonction exponentielle-affine de $r_t$.
\iqlookfor{intégrale gaussienne et forme affine.}
\end{solution}

\begin{solution}{iq:rc:short-rate-models:3}
Construire d'abord l'arbre pour l'état de moyenne nulle~; avancer ensuite pas à pas en propageant les
prix d'Arrow--Debreu de chaque nœud~; à chaque pas, choisir le décalage $\alpha_i$ de sorte que
la somme des prix d'Arrow--Debreu multipliés par les facteurs d'actualisation d'un pas égale le
facteur d'actualisation de la courbe jusqu'à la date suivante~; mettre à jour les prix d'Arrow--Debreu avec
les probabilités de branchement et les taux décalés.
\iqlookfor{récurrence avant avec les prix d'Arrow--Debreu.}
\end{solution}

\begin{solution}{iq:rc:short-rate-models:4}
Les swaptions co-terminales (chaque date d'exercice donnant un swap jusqu'à la
date finale), car la bermudéenne se réduit à l'une d'elles une fois exercée~; un modèle à un facteur
ne peut pas caler toute la matrice, et caler des instruments que le produit
ne fait jamais intervenir gaspille des paramètres sur des risques sans objet.
\iqlookfor{la logique de réplication de l'ensemble de calibration.}
\end{solution}

\begin{solution}{iq:rc:short-rate-models:5}
Produits dépendant d'une corrélation imparfaite entre taux~: options sur spread,
exotiques de type steepener de courbe, et dans une moindre mesure les bermudéennes (valeur de switch). Tester
en valorisant avec un \omterm{def:rc:short-rate-models:g2}{modèle gaussien à deux facteurs} calibré sur les mêmes instruments et
en mesurant l'écart~; s'il est significatif, utiliser le modèle plus riche ou une réserve.
\iqlookfor{corrélation, et un indice de référence pour la mesurer.}
\end{solution}

\begin{solution}{iq:rc:short-rate-models:6}
Vérifier que l'arbre retrouve les facteurs d'actualisation (le calage de $\alpha$)~; diviser le pas par deux et voir
si l'erreur est divisée par deux (discrétisation) ou reste constante (un bogue)~; comparer les options sur obligations zéro-coupon
aux formules fermées~; vérifier l'alignement des dates d'exercice et du calendrier
sur la grille de l'arbre, le branchement aux bords et les probabilités (toutes dans
$[0,1]$, de somme égale à un).
\iqlookfor{étude de convergence et tests unitaires par composant.}
\end{solution}
